\documentclass[12pt, a4paper]{article}

\usepackage[utf8]{inputenc}
\usepackage[T1]{fontenc}
\usepackage{amsmath, amssymb, amsfonts}
\usepackage{amsthm}
\usepackage{mathtools}
\usepackage{tikz}
\usepackage{hyperref}
\usepackage{geometry}
\usepackage{enumitem}
\usepackage{fancyhdr}

\geometry{margin=1in}
\setlength{\headheight}{14.5pt}
\pagestyle{fancy}
\fancyhf{}
\fancyhead[L]{Algebraic Geometry}
\fancyhead[R]{UniMelb 2026 S1}
\fancyfoot[C]{\thepage}

\newtheorem{theorem}{Theorem}
\newtheorem{lemma}[theorem]{Lemma}
\newtheorem{proposition}[theorem]{Proposition}
\newtheorem{corollary}[theorem]{Corollary}
\theoremstyle{definition}
\newtheorem{definition}{Definition}
\newtheorem{example}{Example}
\newtheorem{remark}{Remark}
\newtheorem{exercise}{Exercise}

\DeclareMathOperator{\Spec}{Spec}
\DeclareMathOperator{\GL}{GL}
\DeclareMathOperator{\Hom}{Hom}
\DeclareMathOperator{\coker}{coker}

\title{\textbf{Algebraic Geometry Notes}\\[0.3em]
\large Ch.1: Varieties and Ideals (week 1)}
\author{Geethan Pfeifer}
\date{\today}

\begin{document}

\maketitle

\tableofcontents
\newpage

%\section{Introduction}

%{Overview}
%These notes cover material from the Algebraic Geometry course at University of Melbourne.
\section{Introduction}
This is a slightly reordered, slightly condensed version of the first chapter from \href{https://drive.google.com/file/d/1cJFa9-rGcWDeoQM7paC0mXDbJvZycG4O/view}{Beginning in Algebraic Geometry} by Emily Clader and Dustin Ross. I have excluded examples, proofs, and detailed elaboration; please attend the seminar or read the textbook for these.

\section{Prerequisites}
You should have a rough working idea of the following:
\begin{itemize}
    \item Basic set theory, basic propositional logic, polynomials, functions
    \item Groups, Rings, Fields
    \item Algebraically closed fields
    \item The ring of polynomials, ideals
    \item Basic cartesian geometry
\end{itemize}

\section{Affine Varieties and Vanishing Ideals}
\begin{definition}[Affine Space]
	Given a field $K$, the $n$-dimensional affine space over $K$, denoted $\mathbb{A}^n_K$ is the set of $n$-tuples of $K$.
\end{definition}
Note that the $n$-dimensional vector space over $K$, denoted $K^n$ is \textit{also} the set of $n$-tuples of $K$.  However, for $\mathbb{A}^n_K$ we are not concerned about the vector space structure of $K^n$, rather we treat elements of $\mathbb{A}^n_K$ as inputs to polynomials in $K[x_1, \dots, x_n]$.

\begin{definition}[Polynomial Functions]
	A polynomial $f \in K[x_1, \dots, x_n]$ gives rise to a polynomial function $f : \mathbb{A}_K^n \rightarrow K$, which evaluates the polynomial $f$ for each element in $\mathbb{A}^n_K$. We denote the set of polynomial functions as $K[\mathbb{A}^n]$.
\end{definition}
The map that takes polynomials to polynomial functions is a (ring) isomorphism if and only if $K$ is a infinite field.\footnote{It is also important that $K$ is a field. There are infinite rings that fail to satisfy this property, but we are only concerned about fields here.} Refer to the textbook for a proof.

\begin{definition}[Vanishing Set]
	Given $f \in K[x_1, \dots, x_n]$, we say that $f$ vanishes at $(a_1, \dots, a_n) \in \mathbb{A}^n_K$ if $f$ evaluated at $(a_1, \dots, a_n)$ is $0$. Given a set of polynomials $S \subseteq K[x_1, \dots, x_n]$, the vanishing set of $S$, denoted $\mathcal{V}(S)$, is the set of points in $\mathbb{A}^n_K$ at which \textit{every} polynomial in $S$ vanishes. That is,
	$$\mathcal{V}(S) = \left\{(a_1, \dots, a_n) \in \mathbb{A}^n_K : \forall f \in S, f(a_1, \dots, a_n) = 0\right\}$$
\end{definition}



\begin{definition}[Affine Variety]
	A subset $X \subseteq \mathbb{A}^n_K$ is an affine variety if there exists some subset of polynomials $\mathcal S \subseteq K[x_1, \dots, x_n]$ such that $X = \mathcal{V}(S)$.
\end{definition}
Note that for an affine variety $X$ it may not be possible to choose $S$ with only one member.

\begin{proposition}[Affine Varieties are defined by Ideals]
	For all $S \subseteq K[x_1, \dots, x_n]$, $\mathcal{V}(S) = \mathcal{V}(\langle S \rangle)$.
\end{proposition}
Note that $\langle S \rangle$ is an ideal in the ring $K[x_1, \dots, x_n]$. Refer to the textbook for a proof.

\begin{definition}[Vanishing ideal]
	Let $X \subseteq \mathbb{A}^n_K$. The vanishing ideal of $X$ is defined by:
	$$\mathcal{I}(X) = \{f \in K[x_1, \dots, x_n] : \forall (a_1, \dots, a_n) \in X, f(a_1, \dots, a_n) = 0\}$$
\end{definition}
Note that if $f \in \mathcal{I}(X)$, $(a_1, \dots, a_n) \notin X$ does not imply that $f(a_1, \dots, a_n) \neq 0$. As one might guess from its name, $\mathcal{I}(X)$ is an ideal of $K[x_1, \dots, x_n]$.


\begin{proposition}[Composing $\mathcal{V}$ and $\mathcal{I}$ operators]
	Let $X$ be an affine variety and $S$ a vanishing ideal. Then $\mathcal{V}(\mathcal{I}(X)) = X$ and $\mathcal{I}(\mathcal{V}(S)) = S$.
\end{proposition}
\paragraph{} A slightly generalised version of this proposition is given in the textbook (proposition 1.21).
\paragraph{} We've established a bijection $\mathcal{V}$ between the set of vanishing ideals in $K[x_1, \dots, x_n]$ and the set of affine varieties in in $\mathbb{A}^n_K$ with inverse $\mathcal{I}$.

\begin{definition}[Radical Ideals]
	Let $R$ be a ring, and $I$ an ideal of $R$. $I$ is radical, if it satisfies the following property:
	$$
	\forall a \in R, \forall m \in \mathbb{N}, a^m \in I \implies a \in I
	$$
\end{definition}
In essence this means that radical ideals are closed under radicals (i.e. $n$\textsuperscript{th} roots) that are members of the underlying ring, for all natural $n$.

\begin{definition}[Radical of an ideal]
Let $R$ be a ring, and $I$ an ideal of $R$. Define the radical of $I$ as:
$$
\sqrt I = \left\{
a \in R : \exists m \in \mathbb{N}, a^m \in I
\right\}
$$
\end{definition}
The radical of an ideal is radical, and an ideal $I$ is radical if and only if $I = \sqrt I$.
\begin{proposition}[Vanishing ideals are radical]
	For any set $X \subseteq \mathbb{A}^n_K$, $\mathcal I (X)$ is radical.
\end{proposition}
\begin{proposition}[Radicals of principal ideals]
	Let $R$ be a UFD, $q \in R$. Suppose the distinct irreducible factors of $q$ in $R$ are $q_1, \dots, q_n$. Then, $$
	\sqrt{\langle q \rangle} = \langle q_1 \dots q_n \rangle
	$$
\end{proposition}
\paragraph{}Note that we are taking the ideal generated by the product of $q_1, \dots, q_n$, not the ideal $\langle q_1, \dots, q_n\rangle$.
\paragraph{}This is equivalent to saying that $\langle q \rangle$ is radical iff $q$ is square-free (wrt. the elements of $R$.)

\begin{proposition}[Prime and maximal ideals are radical]
	Every prime ideal is a radical ideal.
\end{proposition}

\begin{theorem}[Nullstellensatz]
   Let $K$ be an algebraically closed field. Let $I \subseteq K[x_1, \ldots, x_n]$ be an ideal. Then $\mathcal I(\mathcal V(I)) = \sqrt{I}$.
\end{theorem}
\begin{corollary}[Bijection between radical ideals and affine varieties]
	If $K$ is an algebraically closed field, then $\mathcal V$ restricted to radical ideals of $K[x_1, \dots x_n]$ is a bijection with the affine varieties of $\mathbb{A}_K^n$ with inverse $\mathcal I$.
\end{corollary}

%\section{References}

%\begin{itemize}
 %   \item Hartshorne, R. \textit{Algebraic Geometry}
%\end{itemize}

\end{document}
